% =====================================================================
\chapter{Ajuste eficiente en parámetros}
\label{tema:4}
% =====================================================================

El ajuste completo de un modelo de $D$ parámetros produce, para cada tarea, una copia de $D$ parámetros y exige guardar durante el entrenamiento el estado del optimizador de todos ellos. Los métodos de \emph{ajuste eficiente en parámetros} (PEFT, del inglés \emph{parameter-efficient fine-tuning}) congelan el modelo preentrenado y entrenan un número pequeño de parámetros nuevos. Este tema presenta los adaptadores, el ajuste de \emph{prompt} y de prefijo, y la noción de dimensión intrínseca que explica por qué bastan tan pocos parámetros. Se sigue el §8.2 de Jurafsky y Martin~\cite{jm3}, la clase~9 de CS224N~\cite{cs224n} y los artículos originales~\cite{houlsby,lester,li-liang,li-intrinsic,aghajanyan}; LoRA, el método más usado, ocupa el \cref{tema:5}.

\section{Planteamiento}

\begin{definicion}[Ajuste eficiente en parámetros]\label{def:peft}
Sea $\theta_0\in\R^{D}$ un modelo preentrenado. Un método PEFT define una familia de modelos $f_{\theta_0,\varphi}$ con parámetros nuevos $\varphi\in\R^{P}$, $P\ll D$, tal que existe $\varphi_0$ con $f_{\theta_0,\varphi_0}=f_{\theta_0}$, y entrena solo $\varphi$, partiendo de $\varphi_0$ o de un valor próximo, con $\theta_0$ congelado.
\end{definicion}

Los métodos se agrupan en tres familias: \emph{aditivos}, que añaden módulos o vectores nuevos (adaptadores, ajuste de \emph{prompt} y de prefijo); \emph{selectivos}, que entrenan un subconjunto de los parámetros existentes (la sonda lineal del \cref{tema:1} es el caso extremo), y de \emph{reparametrización}, que escriben la actualización $\Delta\theta$ como función de pocos parámetros (subespacios aleatorios, \cref{sec:intrinseca}, y LoRA, \cref{tema:5}).

\begin{proposicion}[Memoria del entrenamiento]\label{prop:memoria-peft}
Con precisión mixta y Adam, el ajuste completo requiere unos $16D$ bytes para pesos, gradientes y estado del optimizador (módulo~17). Con un método PEFT y el modelo base en bf16, la memoria es de unos $2D+16P$ bytes, sin contar las activaciones.
\end{proposicion}

Con $D=7\cdot 10^{9}$ y $P=2\cdot 10^{7}$, la memoria baja de $112$~GB a $14+0{,}32=14{,}32$~GB. Además, cada tarea se almacena con $P$ parámetros en lugar de $D$, y un mismo modelo base puede servir muchas tareas cambiando solo $\varphi$.

\section{Adaptadores}

\begin{definicion}[Adaptador]\label{def:adaptador}
Sea $h\in\R^{d}$ la salida de una subcapa de un Transformer. Un \emph{adaptador} con dimensión de cuello $m\ll d$ calcula
\[
a(h)=h+W_{\uparrow}\,g\bigl(W_{\downarrow}h+b_{\downarrow}\bigr)+b_{\uparrow},\qquad W_{\downarrow}\in\R^{m\times d},\ W_{\uparrow}\in\R^{d\times m},
\]
con $g$ una no linealidad. Houlsby \emph{et al.}~\cite{houlsby} insertan dos adaptadores por capa, tras la subcapa de atención y tras la red de alimentación hacia delante, y entrenan además los parámetros de normalización de capa.
\end{definicion}

\addtocounter{figure}{1}%

\begin{proposicion}\label{prop:adaptador}
Un adaptador tiene $2md+m+d$ parámetros. Si se inicializa con $W_{\uparrow}=0$ y $b_{\uparrow}=0$, entonces $a(h)=h$ para todo $h$, de modo que el modelo con adaptadores coincide al principio con el preentrenado.
\end{proposicion}

Houlsby \emph{et al.} usan una inicialización \emph{casi} nula, de modo que el adaptador es casi la identidad, y obtienen en GLUE resultados a menos de un $0{,}4\,\%$ de los del ajuste completo añadiendo un $3{,}6\,\%$ de parámetros por tarea~\cite{houlsby}.

\begin{observacion}[Latencia]
Los adaptadores añaden capas que se ejecutan en serie y no pueden fusionarse con los pesos originales, porque contienen una no linealidad. Hu \emph{et al.}~\cite[§3]{hu-lora} midieron aumentos de la latencia de inferencia superiores al $30\,\%$ con secuencias cortas y lotes pequeños. LoRA evita este problema (\cref{tema:5}).
\end{observacion}

\section{Ajuste de \emph{prompt} y de prefijo}

\begin{definicion}[Ajuste de \emph{prompt}]\label{def:prompt-tuning}
Sea $E\in\R^{T\times d}$ la matriz de \emph{embeddings} de una entrada. El \emph{ajuste de prompt} antepone $p$ vectores entrenables $P_e\in\R^{p\times d}$ y alimenta el modelo congelado con $[P_e;E]\in\R^{(p+T)\times d}$. Tiene $pd$ parámetros~\cite{lester}.
\end{definicion}

Los vectores de $P_e$ no tienen por qué corresponder a ningún token del vocabulario: forman un \emph{prompt continuo} (\emph{soft prompt}). Lester \emph{et al.}~\cite{lester} observaron que la diferencia con el ajuste completo disminuye con la escala y prácticamente desaparece en modelos de unos $10^{10}$ parámetros.

\begin{definicion}[Ajuste de prefijo]\label{def:prefix}
El \emph{ajuste de prefijo} antepone, en cada una de las $L$ capas, $p$ claves y $p$ valores entrenables, $P_K^{l},P_V^{l}\in\R^{p\times d}$, a las claves y valores calculados a partir de la entrada. Tiene $2Lpd$ parámetros. Para estabilizar el entrenamiento, Li y Liang~\cite{li-liang} reparametrizan los prefijos mediante una pequeña red que se descarta al terminar; con ello entrenan en torno al $0{,}1\,\%$ de los parámetros.
\end{definicion}

El efecto de un prefijo sobre la atención admite una descripción exacta.

\begin{teorema}[El prefijo interpola la atención]\label{teo:prefijo}
Sean $q\in\R^{d}$ una consulta, $K,V\in\R^{T\times d}$ las claves y valores de la entrada y $P_K,P_V\in\R^{p\times d}$ los del prefijo. Escríbase $\mathrm{At}(q,K,V)=\softmax(Kq/\sqrt{d})\T V$. Entonces
\[
\mathrm{At}\bigl(q,[P_K;K],[P_V;V]\bigr)=\bigl(1-\mu(q)\bigr)\,\mathrm{At}(q,K,V)+\mu(q)\,\mathrm{At}(q,P_K,P_V),
\]
con
\[
\mu(q)=\frac{S_P}{S_P+S_K},\qquad S_P=\sum_{j=1}^{p}e^{(P_K)_j\cdot q/\sqrt{d}},\quad S_K=\sum_{i=1}^{T}e^{k_i\cdot q/\sqrt{d}}.
\]
\end{teorema}

\begin{corolario}\label{cor:prefijo}
El ajuste de prefijo no altera los pesos relativos de atención entre las posiciones de la entrada: solo mezcla la salida original con un vector que depende de la consulta, con una proporción $\mu(q)\in(0,1)$.
\end{corolario}

Esta lectura, que relaciona el ajuste de prefijo con los adaptadores, se desarrolla en He \emph{et al.}~\cite{he-unified}.

\section{Dimensión intrínseca}\label{sec:intrinseca}

¿Por qué basta con ajustar una fracción mínima de los parámetros? Li \emph{et al.}~\cite{li-intrinsic} propusieron medir la dimensión efectiva de un problema de optimización restringiendo el entrenamiento a un subespacio aleatorio.

\begin{definicion}[Subespacio aleatorio y dimensión intrínseca]\label{def:intrinseca}
Sean $\theta_0\in\R^{D}$ y $P\in\R^{D\times d'}$ una matriz aleatoria fija. Se entrena solo $z\in\R^{d'}$, inicializado en $0$, con la reparametrización $\theta=\theta_0+Pz$. La \emph{dimensión intrínseca} $d_{90}$ de una tarea es el menor $d'$ con el que se alcanza el $90\,\%$ del rendimiento del ajuste completo~\cite[§3]{aghajanyan}.
\end{definicion}

\begin{proposicion}[Gradiente en el subespacio]\label{prop:grad-subespacio}
Si $\mathcal{L}$ es derivable, $\nabla_z\mathcal{L}(\theta_0+Pz)=P\T\,\nabla_\theta\mathcal{L}(\theta)\big|_{\theta=\theta_0+Pz}$.
\end{proposicion}

Aghajanyan \emph{et al.}~\cite{aghajanyan} estimaron la dimensión intrínseca del ajuste de modelos de lenguaje preentrenados, usando para $P$ la transformada Fastfood, que evita almacenar una matriz densa de $D\times d'$. Sus tres resultados principales son:
\begin{enumerate}
  \item con RoBERTa-Large, unos 200 parámetros bastan para alcanzar el $90\,\%$ del rendimiento del ajuste completo en la tarea MRPC~\cite[§4.2]{aghajanyan};
  \item la dimensión intrínseca disminuye a lo largo del preentrenamiento y es menor en los modelos con más parámetros~\cite[§5.1 y~§5.2]{aghajanyan};
  \item una cota de generalización basada en compresión depende de la dimensión intrínseca y no del número total de parámetros: con $n$ ejemplos, el error de generalización es $O(\sqrt{d_{90}/n})$, para modelos entrenados en el subespacio~\cite[§5.3]{aghajanyan}.
\end{enumerate}

\begin{observacion}
El segundo resultado explica en parte la paradoja de que los modelos mayores se ajusten con menos datos: el preentrenamiento comprime el problema de adaptación. Hu \emph{et al.} formularon a partir de aquí la hipótesis que motiva LoRA~\cite[§1 y~§4.1]{hu-lora}: si el ajuste vive en un subespacio de dimensión pequeña, también la actualización de cada matriz de pesos debe tener un \emph{rango intrínseco} bajo. LoRA no fija las direcciones al azar: aprende la mejor base de rango $r$.
\end{observacion}

